\chapter{2020年全国II卷（理）}

\section{单选题}

\begin{problem}
已知集合 \(U=\{-2,-1,0,1,2,3\}\)，\(A=\{-1,0,1\}\)，\(B=\{1,2\}\)，则 \(\complement_U\left( A\cup B \right)=\)（\quad）

\choices
    {\(\{-2,3\}\)}
    {\(\{-2,2,3\}\)}
    {\(\{-2,-1,0,3\}\)}
    {\(\{-2,-1,0,2,3\}\)}

\end{problem}

\begin{answer}
A
\end{answer}

\begin{solution}
集合 \(A\cup B=\{-1,0,1,2\}\). 所以 \(\complement_U\left( A\cup B \right)=\{-2,3\}\). 故选 \(A\).
\end{solution}

\begin{problem}
若 \(\alpha\) 为第四象限角，则（\quad）

\choices
    {\(\cos 2\alpha>0\)}
    {\(\cos 2\alpha<0\)}
    {\(\sin 2\alpha>0\)}
    {\(\sin 2\alpha<0\)}

\end{problem}

\begin{answer}
D
\end{answer}

\begin{solution}
因为 \(\alpha\) 为第四象限角，所以 \(-\frac{\pi}{2}+2k\pi<\alpha<2k\pi\)，\(k\in\Z\). 于是 \(-\pi+4k\pi<2\alpha<4k\pi\). 所以 \(2\alpha\) 是第三或第四象限角，或为 \(y\) 轴负半轴上的角，从而 \(\sin 2\alpha<0\). 故选 \(D\).
\end{solution}

\begin{problem}
在新冠肺炎疫情防控期间，某超市开通网上销售业务，每天能完成 \(1200\) 份订单的配货，由于订单量大幅增加，导致订单积压. 为解决困难，许多志愿者踊跃报名参加配货工作. 已知该超市某日积压 \(500\) 份订单未配货，预计第二天的新订单超过 \(1600\) 份的概率为 \(0.05\). 志愿者每人每天能完成 \(50\) 份订单的配货，为使第二天完成积压订单及当日订单的配货的概率不小于 \(0.95\)，则至少需要志愿者（\quad）

\choices
    {\(10\) 名}
    {\(18\) 名}
    {\(24\) 名}
    {\(32\) 名}

\end{problem}

\begin{answer}
B
\end{answer}

\begin{solution}
第二天的新订单不超过 \(1600\) 份的概率为 \(0.95\). 要使第二天完成积压订单及当日订单的配货的概率不小于 \(0.95\)，只需保证在新订单为 \(1600\) 份时仍能完成配货. 此时共需完成 \(500+1600=2100\) 份订单，超市本身能完成 \(1200\) 份，还需完成 \(2100-1200=900\) 份. 因为每名志愿者每天能完成 \(50\) 份订单，所以至少需要 \(\frac{900}{50}=18\) 名志愿者. 故选 \(B\).
\end{solution}

\begin{problem}
北京天坛的圜丘坛为古代祭天的场所，分上，中，下三层. 上层中心有一块圆形石板（称为天心石），环绕天心石砌 \(9\) 块扇面形石板构成第一环，向外每环依次增加 \(9\) 块. 下一层的第一环比上一层的最后一环多 \(9\) 块，向外每环依次也增加 \(9\) 块. 已知每层环数相同，且下层比中层多 \(729\) 块，则三层共有扇面形石板（不含天心石）（\quad）

\choices
    {\(3699\) 块}
    {\(3474\) 块}
    {\(3402\) 块}
    {\(3339\) 块}

\end{problem}

\begin{answer}
C
\end{answer}

\begin{solution}
设每层有 \(n\) 环. 由题意可知，从内到外每环石板数构成公差为 \(9\) 的等差数列，且首项为 \(9\). 记前三层对应等差数列前 \(n\) 项和分别为 \(S_n\)，\(S_{2n}-S_n\)，\(S_{3n}-S_{2n}\). 因为下层比中层多 \(729\) 块，所以 \(\left( S_{3n}-S_{2n} \right)-\left( S_{2n}-S_n \right)=729\). 由等差数列性质可得 \(\left( S_{3n}-S_{2n} \right)-\left( S_{2n}-S_n \right)=n^2\cdot 9\). 于是 \(9n^2=729\)，解得 \(n=9\). 所以三层共有扇面形石板 \(S_{27}=27\times 9+\frac{27\times 26}{2}\times 9=3402\). 故选 \(C\).
\end{solution}

\begin{problem}
若过点 \((2,1)\) 的圆与两坐标轴都相切，则圆心到直线 \(2x-y-3=0\) 的距离为（\quad）

\choices
    {\(\dfrac{\sqrt5}{5}\)}
    {\(\dfrac{2\sqrt5}{5}\)}
    {\(\dfrac{3\sqrt5}{5}\)}
    {\(\dfrac{4\sqrt5}{5}\)}

\end{problem}

\begin{answer}
B
\end{answer}

\begin{solution}
由题意可知所求圆在第一象限. 设圆心为 \((a,a)\)，半径为 \(a\)（\(a>0\)），则圆的方程为 \(\left( x-a \right)^2+\left( y-a \right)^2=a^2\). 把点 \((2,1)\) 代入，得 \(\left( 2-a \right)^2+\left( 1-a \right)^2=a^2\)，解得 \(a=1\) 或 \(a=5\).

所以圆心为 \((1,1)\) 或 \((5,5)\). 圆心到直线 \(2x-y-3=0\) 的距离分别为 \(\frac{|2-1-3|}{\sqrt5}=\frac{2}{\sqrt5}\) 和 \(\frac{|10-5-3|}{\sqrt5}=\frac{2}{\sqrt5}\). 故所求距离为 \(\frac{2}{\sqrt5}=\frac{2\sqrt5}{5}\). 故选 \(B\).
\end{solution}

\begin{problem}
数列 \(\{a_n\}\) 中，\(a_1=2\)，\(a_{m+n}=a_m a_n\). 若 \(a_{k+1}+a_{k+2}+\cdots+a_{k+10}=2^{15}-2^5\)，则 \(k=\)（\quad）

\choices
    {\(2\)}
    {\(3\)}
    {\(4\)}
    {\(5\)}

\end{problem}

\begin{answer}
C
\end{answer}

\begin{solution}
由 \(a_1=2\) 且 \(a_{m+n}=a_m a_n\)，取 \(m=1\)，得 \(a_{n+1}=a_1a_n=2a_n\). 所以数列 \(\{a_n\}\) 是首项为 \(2\)，公比为 \(2\) 的等比数列，即 \(a_n=2^n\). 于是 \(a_{k+1}+a_{k+2}+\cdots+a_{k+10}=2^{k+1}+2^{k+2}+\cdots+2^{k+10}=2^{k+1}\left( 2^{10}-1 \right)\). 而 \(2^{15}-2^5=2^5\left( 2^{10}-1 \right)\). 所以 \(k+1=5\)，得 \(k=4\). 故选 \(C\).
\end{solution}

\begin{problem}
如图是一个多面体的三视图，这个多面体某条棱的一个端点在正视图中对应的点为 \(M\)，在俯视图中对应的点为 \(N\)，则该端点在侧视图中对应的点为（\quad）

\begin{center}
\bitmapinclude[max width=6.0cm,max totalheight=2.4cm]{img/2020/national_paper_2_science/q7_fig1.png}
\end{center}

\choices
    {\(E\)}
    {\(F\)}
    {\(G\)}
    {\(H\)}

\end{problem}

\begin{answer}
A
\end{answer}

\begin{solution}
根据三视图还原直观图，并由三视图与直观图中点的对应关系可知，这个端点在侧视图中对应的点为 \(E\). 故选 \(A\).
\end{solution}

\begin{problem}
设 \(O\) 为坐标原点，直线 \(x=a\) 与双曲线 \(C:\dfrac{x^2}{a^2}-\dfrac{y^2}{b^2}=1\)（\(a>0\)，\(b>0\)）的两条渐近线分别交于 \(D\)，\(E\) 两点. 若 \(\triangle ODE\) 的面积为 \(8\)，则 \(C\) 的焦距的最小值为（\quad）

\choices
    {\(4\)}
    {\(8\)}
    {\(16\)}
    {\(32\)}

\end{problem}

\begin{answer}
B
\end{answer}

\begin{solution}
双曲线 \(C\) 的渐近线方程为 \(y=\pm \frac{b}{a}x\). 直线 \(x=a\) 与两条渐近线交于 \(D(a,b)\)，\(E(a,-b)\). 所以 \(S_{\triangle ODE}=\frac{1}{2}\cdot 2b\cdot a=ab=8\). 双曲线的焦距为 \(2\sqrt{a^2+b^2}\). 由 \(a^2+b^2\ge 2ab=16\) 可得 \(2\sqrt{a^2+b^2}\ge 2\sqrt{16}=8\). 当 \(a=b\) 时取等号，所以焦距的最小值为 \(8\). 故选 \(B\).
\end{solution}

\begin{problem}
设函数 \(f(x)=\ln |2x+1|-\ln |2x-1|\)，则 \(f(x)\)（\quad）

\choices
    {是偶函数，且在 \(\left( \dfrac{1}{2},+\infty \right)\) 单调递增}
    {是奇函数，且在 \(\left( -\dfrac{1}{2},\dfrac{1}{2} \right)\) 单调递减}
    {是偶函数，且在 \(\left( -\infty,-\dfrac{1}{2} \right)\) 单调递增}
    {是奇函数，且在 \(\left( -\infty,-\dfrac{1}{2} \right)\) 单调递减}

\end{problem}

\begin{answer}
D
\end{answer}

\begin{solution}
因为 \(f(-x)=\ln |1-2x|-\ln |-(2x+1)|=\ln |2x-1|-\ln |2x+1|=-f(x)\)，所以 \(f(x)\) 是奇函数. 又 \(f'(x)=\frac{2}{2x+1}-\frac{2}{2x-1}=-\frac{4}{4x^2-1}\). 当 \(x\in\left( -\infty,-\dfrac{1}{2} \right)\) 时，\(4x^2-1>0\)，故 \(f'(x)<0\)，所以 \(f(x)\) 在 \(\left( -\infty,-\dfrac{1}{2} \right)\) 上单调递减. 故选 \(D\).
\end{solution}

\begin{problem}
已知 \(\triangle ABC\) 是面积为 \(\dfrac{9\sqrt3}{4}\) 的等边三角形，且其顶点都在球 \(O\) 的球面上. 若球 \(O\) 的表面积为 \(16\pi\)，则 \(O\) 到平面 \(ABC\) 的距离为（\quad）

\choices
    {\(\sqrt3\)}
    {\(\dfrac{3}{2}\)}
    {\(1\)}
    {\(\dfrac{\sqrt3}{2}\)}

\end{problem}

\begin{answer}
C
\end{answer}

\begin{solution}
设等边三角形边长为 \(s\)，由面积公式得 \(\frac{\sqrt3}{4}s^2=\frac{9\sqrt3}{4}\)，解得 \(s=3\). 所以 \(\triangle ABC\) 的外接圆半径为 \(R_1=\frac{s}{\sqrt3}=\sqrt3\). 又球 \(O\) 的表面积为 \(16\pi\)，故球半径 \(R\) 满足 \(4\pi R^2=16\pi\)，解得 \(R=2\). 因此球心 \(O\) 到平面 \(ABC\) 的距离为 \(\sqrt{R^2-R_1^2}=\sqrt{4-3}=1\). 故选 \(C\).
\end{solution}

\begin{problem}
若 \(2^x-2^y<3^{-x}-3^{-y}\)，则（\quad）

\choices
    {\(\ln \left( y-x+1 \right)>0\)}
    {\(\ln \left( y-x+1 \right)<0\)}
    {\(\ln \left| x-y \right|>0\)}
    {\(\ln \left| x-y \right|<0\)}

\end{problem}

\begin{answer}
A
\end{answer}

\begin{solution}
由 \(2^x-2^y<3^{-x}-3^{-y}\) 可得 \(2^x-3^{-x}<2^y-3^{-y}\). 令 \(f(x)=2^x-3^{-x}\)，则 \(f'(x)=2^x\ln 2+3^{-x}\ln 3>0\). 所以 \(f(x)\) 在 \(\R\) 上单调递增. 由 \(f(x)<f(y)\) 得 \(x<y\)，即 \(y-x>0\)，于是 \(y-x+1>1\). 故 \(\ln \left( y-x+1 \right)>\ln 1=0\). 故选 \(A\).
\end{solution}

\begin{problem}
\(0\)-\(1\) 周期序列在通信技术中有着重要应用. 若序列 \(a_1a_2\cdots a_n\cdots\) 满足 \(a_i\in\{0,1\}\)（\(i=1,2,\cdots\)），且存在正整数 \(m\)，使得 \(a_{i+m}=a_i\)（\(i=1,2,\cdots\)）成立，则称其为 \(0\)-\(1\) 周期序列，并称满足 \(a_{i+m}=a_i\)（\(i=1,2,\cdots\)）的最小正整数 \(m\) 为这个序列的周期. 对于周期为 \(m\) 的 \(0\)-\(1\) 序列 \(a_1a_2\cdots a_n\cdots\)，\(C(k)=\dfrac{1}{m}\sum_{i=1}^{m} a_i a_{i+k}\)（\(k=1,2,\cdots,m-1\)）是描述其性质的重要指标，下列周期为 \(5\) 的 \(0\)-\(1\) 序列中，满足 \(C(k)\le \dfrac{1}{5}\)（\(k=1,2,3,4\)）的序列是（\quad）

\choices
    {\(11010\cdots\)}
    {\(11011\cdots\)}
    {\(10001\cdots\)}
    {\(11001\cdots\)}

\end{problem}

\begin{answer}
C
\end{answer}

\begin{solution}
分别对四个选项中 \(k=1,2,3,4\) 进行讨论.

对于 \(A\) 选项，序列为 \(11010\ 11010\)，则 \(C(1)=\frac{1}{5}(1+0+0+0+0)=\frac{1}{5}\)，\(C(2)=\frac{1}{5}(0+1+0+1+0)=\frac{2}{5}>\frac{1}{5}\)，不满足条件，排除.

对于 \(B\) 选项，序列为 \(11011\ 11011\)，则 \(C(1)=\frac{1}{5}(1+0+0+1+1)=\frac{3}{5}>\frac{1}{5}\)，不满足条件，排除.

对于 \(C\) 选项，序列为 \(10001\ 10001\ 10001\)，则 \(C(1)=\frac{1}{5}(0+0+0+0+1)=\frac{1}{5}\)，\(C(2)=\frac{1}{5}(0+0+0+0+0)=0\)，\(C(3)=\frac{1}{5}(0+0+0+0+0)=0\)，\(C(4)=\frac{1}{5}(1+0+0+0+0)=\frac{1}{5}\)，均满足条件.

对于 \(D\) 选项，序列为 \(11001\ 11001\)，则 \(C(1)=\frac{1}{5}(1+0+0+0+1)=\frac{2}{5}>\frac{1}{5}\)，不满足条件.

故选 \(C\).
\end{solution}
\section{填空题}

\begin{problem}
已知单位向量 \(\bs{a},\bs{b}\) 的夹角为 \(45^\circ\)，\(k\bs{a}-\bs{b}\) 与 \(\bs{a}\) 垂直，则 \(k=\fillinblank{}\)

\end{problem}

\begin{answer}
\(\dfrac{\sqrt2}{2}\)
\end{answer}

\begin{solution}
因为 \(\bs{a},\bs{b}\) 为单位向量，且 \(\bs{a},\bs{b}\) 的夹角为 \(45^\circ\)，所以 \(\bs{a}\cdot\bs{b}=|\bs{a}|\cdot|\bs{b}|\cos 45^\circ=1\times 1\times \frac{\sqrt2}{2}=\frac{\sqrt2}{2}\). 又因为 \(k\bs{a}-\bs{b}\) 与 \(\bs{a}\) 垂直，所以 \(\left( k\bs{a}-\bs{b} \right)\cdot \bs{a}=k|\bs{a}|^2-\bs{a}\cdot\bs{b}=0\). 即 \(k-\frac{\sqrt2}{2}=0\)，所以 \(k=\frac{\sqrt2}{2}\). 故答案为 \(\dfrac{\sqrt2}{2}\).
\end{solution}

\begin{problem}
\(4\) 名同学到 \(3\) 个小区参加垃圾分类宣传活动，每名同学只去 \(1\) 个小区，每个小区至少安排 \(1\) 名同学，则不同的安排方法共有\fillinblank{} 种

\end{problem}

\begin{answer}
\(36\)
\end{answer}

\begin{solution}
因为有一个小区安排 \(2\) 人，其余两个小区各安排 \(1\) 人，所以不同的安排方法共有 \(C_4^2A_3^3=36\) 种. 故答案为 \(36\).
\end{solution}

\begin{problem}
设复数 \(z_1,z_2\) 满足 \(|z_1|=|z_2|=2\)，\(z_1+z_2=\sqrt3+\i\)，则 \(|z_1-z_2|=\fillinblank{}\)

\end{problem}

\begin{answer}
\(2\sqrt3\)
\end{answer}

\begin{solution}
因为 \(|z_1+z_2|=|\sqrt3+\i|=2\)，所以 \(|z_1+z_2|^2=\left( z_1+z_2 \right)\overline{\left( z_1+z_2 \right)}=4\). 即 \(|z_1|^2+|z_2|^2+z_1\overline{z_2}+\overline{z_1}z_2=4\). 由 \(|z_1|=|z_2|=2\) 得 \(8+z_1\overline{z_2}+\overline{z_1}z_2=4\)，所以 \(z_1\overline{z_2}+\overline{z_1}z_2=-4\). 于是 \(|z_1-z_2|^2=\left( z_1-z_2 \right)\overline{\left( z_1-z_2 \right)}=|z_1|^2+|z_2|^2-\left( z_1\overline{z_2}+\overline{z_1}z_2 \right)=12\). 又 \(|z_1-z_2|>0\)，故 \(|z_1-z_2|=2\sqrt3\). 故答案为 \(2\sqrt3\).
\end{solution}

\begin{problem}
设有下列四个命题：

\begin{enumerate}
\item \(p_1\)：两两相交且不过同一点的三条直线必在同一平面内；
\item \(p_2\)：过空间中任意三点有且仅有一个平面；
\item \(p_3\)：若空间两条直线不相交，则这两条直线平行；
\item \(p_4\)：若直线 \(l\subset \alpha\)，直线 \(m\perp \alpha\)，则 \(m\perp l\).
\end{enumerate}

则下述命题中所有真命题的序号是\fillinblank{}

\begin{enumerate}
\item \(\circled{1}\ p_1\land p_4\)
\item \(\circled{2}\ p_1\land p_2\)
\item \(\circled{3}\ \neg p_2\lor p_3\)
\item \(\circled{4}\ \neg p_3\lor \neg p_4\)
\end{enumerate}

\end{problem}

\begin{answer}
\(\circled{1}\circled{3}\circled{4}\)
\end{answer}

\begin{solution}
命题 \(p_1\) 为真命题，因为两两相交且不过同一点的三条直线必在同一平面内.

命题 \(p_2\) 为假命题，因为当三点共线时，过这三点的平面有无数个.

命题 \(p_3\) 为假命题，因为空间中两条直线不相交，也可能是异面直线.

命题 \(p_4\) 为真命题，因为若直线 \(m\perp \alpha\)，则 \(m\) 垂直于平面 \(\alpha\) 内任意一条过垂足的直线，从而 \(m\perp l\).

所以复合命题中，\(\circled{1}\ p_1\land p_4\) 为真命题，\(\circled{2}\ p_1\land p_2\) 为假命题，\(\circled{3}\ \neg p_2\lor p_3\) 为真命题，\(\circled{4}\ \neg p_3\lor \neg p_4\) 为真命题.

故真命题的序号是 \(\circled{1}\circled{3}\circled{4}\).
\end{solution}
\section{解答题}

\begin{problem}
\(\triangle ABC\) 中，\(\sin^2 A-\sin^2 B-\sin^2 C=\sin B\sin C\).

\begin{enumerate}
\item 求 \(A\)；
\item 若 \(BC=3\)，求 \(\triangle ABC\) 周长的最大值.
\end{enumerate}

\end{problem}

\begin{solution}
（1）设 \(\triangle ABC\) 的内角 \(A,B,C\) 所对的边分别为 \(a,b,c\). 因为 \(\sin^2 A-\sin^2 B-\sin^2 C=\sin B\sin C\)，由正弦定理可得 \(a^2-b^2-c^2=bc\)，即 \(b^2+c^2-a^2=-bc\). 由余弦定理可得 \(\cos A=\frac{b^2+c^2-a^2}{2bc}=-\frac{bc}{2bc}=-\frac{1}{2}\). 由 \(0<A<\pi\)，可得 \(A=\frac{2\pi}{3}\).

（2）由题意可得 \(a=3\). 又 \(B+C=\frac{\pi}{3}\)，可设 \(B=\frac{\pi}{6}-d\)，\(C=\frac{\pi}{6}+d\)，\(-\frac{\pi}{6}<d<\frac{\pi}{6}\). 由正弦定理可得 \(\frac{3}{\sin \frac{2\pi}{3}}=\frac{b}{\sin B}=\frac{c}{\sin C}=2\sqrt3\)，所以 \(b=2\sqrt3\sin \left( \frac{\pi}{6}-d \right)\)，\(c=2\sqrt3\sin \left( \frac{\pi}{6}+d \right)\).

于是 \(\triangle ABC\) 的周长为
\examdisplayaligned{
a+b+c
&=3+2\sqrt3\left[ \sin \left( \frac{\pi}{6}-d \right)+\sin \left( \frac{\pi}{6}+d \right) \right] \\
&=3+2\sqrt3\left( \frac{1}{2}\cos d-\frac{\sqrt3}{2}\sin d+\frac{1}{2}\cos d+\frac{\sqrt3}{2}\sin d \right) \\
&=3+2\sqrt3\cos d.
}

当 \(d=0\)，即 \(B=C=\frac{\pi}{6}\) 时，\(\triangle ABC\) 的周长取得最大值 \(3+2\sqrt3\).
\end{solution}

\begin{problem}
某沙漠地区经过治理，生态系统得到很大改善，野生动物数量有所增加. 为调查该地区某种野生动物的数量，将其分成面积相近的 \(200\) 个地块，从这些地块中用简单随机抽样的方法抽取 \(20\) 个作为样区，调查得到样本数据 \((x_i,y_i)\)（\(i=1,2,\cdots,20\)），其中 \(x_i\) 和 \(y_i\) 分别表示第 \(i\) 个样区的植物覆盖面积（单位：公顷）和这种野生动物的数量，并计算得 \(\sum_{i=1}^{20}x_i=60\)，\(\sum_{i=1}^{20}y_i=1200\)，\(\sum_{i=1}^{20}\left( x_i-\bar{x} \right)^2=80\)，\(\sum_{i=1}^{20}\left( y_i-\bar{y} \right)^2=9000\)，\(\sum_{i=1}^{20}\left( x_i-\bar{x} \right)\left( y_i-\bar{y} \right)=800\).

\begin{enumerate}
\item 求该地区这种野生动物数量的估计值（这种野生动物数量的估计值等于样区这种野生动物数量的平均数乘以地块数）；
\item 求样本 \((x_i,y_i)\)（\(i=1,2,\cdots,20\)）的相关系数（精确到 \(0.01\)）；
\item 根据现有统计资料，各地块间植物覆盖面积差异很大. 为提高样本的代表性以获得该地区这种野生动物数量更准确的估计，请给出一种你认为更合理的抽样方法，并说明理由.
\end{enumerate}

附：相关系数 \(r=\dfrac{\sum_{i=1}^{n}\left( x_i-\bar{x} \right)\left( y_i-\bar{y} \right)}{\sqrt{\sum_{i=1}^{n}\left( x_i-\bar{x} \right)^2\sum_{i=1}^{n}\left( y_i-\bar{y} \right)^2}}\)，\(\sqrt2\approx 1.414\).

\end{problem}

\begin{solution}
（1）由已知，\(\bar{y}=\frac{1}{20}\sum_{i=1}^{20}y_i=\frac{1200}{20}=60\). 所以该地区这种野生动物数量的估计值为 \(60\times 200=12000\).

（2）由已知，\(\sum_{i=1}^{20}\left( x_i-\bar{x} \right)^2=80\)，\(\sum_{i=1}^{20}\left( y_i-\bar{y} \right)^2=9000\)，\(\sum_{i=1}^{20}\left( x_i-\bar{x} \right)\left( y_i-\bar{y} \right)=800\). 所以 \(r=\frac{800}{\sqrt{80\times 9000}}=\frac{800}{600\sqrt2}=\frac{2\sqrt2}{3}\approx 0.94\).

（3）更合理的抽样方法是分层抽样. 根据植物覆盖面积的大小对地块分层，再对 \(200\) 个地块进行分层抽样.

理由如下：由（2）知各样区的这种野生动物数量与植物覆盖面积有很强的正相关. 由于各地块间植物覆盖面积差异很大，从而各地块间这种野生动物数量差异也很大，采用分层抽样的方法较好地保持了样本结构与总体结构的一致性，提高了样本的代表性，从而可以获得该地区这种野生动物数量更准确的估计.
\end{solution}

\begin{problem}
已知椭圆 \(C_1:\dfrac{x^2}{a^2}+\dfrac{y^2}{b^2}=1\)（\(a>b>0\)）的右焦点 \(F\) 与抛物线 \(C_2\) 的焦点重合，\(C_1\) 的中心与 \(C_2\) 的顶点重合，过 \(F\) 且与 \(x\) 轴垂直的直线交 \(C_1\) 于 \(A,B\) 两点，交 \(C_2\) 于 \(C,D\) 两点，且 \(|CD|=\dfrac{4}{3}|AB|\).

\begin{enumerate}
\item 求 \(C_1\) 的离心率；
\item 设 \(M\) 是 \(C_1\) 与 \(C_2\) 的公共点. 若 \(|MF|=5\)，求 \(C_1\) 与 \(C_2\) 的标准方程.
\end{enumerate}

\end{problem}

\begin{solution}
（1）因为 \(F\) 为 \(C_1\) 的焦点，且 \(AB\perp x\) 轴，所以 \(F(c,0)\)，\(|AB|=\frac{2b^2}{a}\). 设 \(C_2\) 的标准方程为 \(y^2=2px\)（\(p>0\)）. 因为 \(F\) 为 \(C_2\) 的焦点，且 \(CD\perp x\) 轴，所以 \(F\left( \frac{p}{2},0 \right)\)，\(|CD|=2p\). 由 \(|CD|=\dfrac{4}{3}|AB|\)，且 \(C_1,C_2\) 的焦点重合，得 \(\begin{cases}
c=\dfrac{p}{2},\\[4pt]
2p=\dfrac{4}{3}\cdot \dfrac{2b^2}{a}.
\end{cases}\). 消去 \(p\)，可得 \(4c=\frac{8b^2}{3a}\). 所以 \(3ac=2b^2\). 又因为 \(b^2=a^2-c^2\)，所以 \(3ac=2a^2-2c^2\). 设 \(C_1\) 的离心率为 \(e\)，则 \(e=\dfrac{c}{a}\). 上式化为 \(2e^2+3e-2=0\). 解得 \(e=\frac{1}{2}\)（\(-2\text{ 舍去}\)）. 故 \(C_1\) 的离心率为 \(\dfrac{1}{2}\).

（2）由（1）可得 \(a=2c\)，\(b=\sqrt3\,c\)，\(p=2c\). 所以 \(C_1:\frac{x^2}{4c^2}+\frac{y^2}{3c^2}=1\)，\(C_2:y^2=4cx\). 联立两曲线方程，消去 \(y\)，可得 \(3x^2+16cx-12c^2=0\). 所以 \(\left( 3x-2c \right)\left( x+6c \right)=0\). 解得 \(x=\frac{2}{3}c\) 或 \(x=-6c\). 后者舍去，所以 \(x=\frac{2}{3}c\). 由抛物线的定义，\(|MF|=x+\frac{p}{2}=\frac{2}{3}c+c=\frac{5}{3}c=5\). 解得 \(c=3\). 所以 \(C_1:\frac{x^2}{36}+\frac{y^2}{27}=1\)，\(C_2:y^2=12x\).
\end{solution}

\begin{problem}
如图，已知三棱柱 \(ABC-A_1B_1C_1\) 的底面是正三角形，侧面 \(BB_1C_1C\) 是矩形，\(M,N\) 分别为 \(BC,B_1C_1\) 的中点，\(P\) 为 \(AM\) 上一点. 过 \(B_1C_1\) 和 \(P\) 的平面交 \(AB\) 于 \(E\)，交 \(AC\) 于 \(F\).

\begin{enumerate}
\item 证明：\(AA_1//MN\)，且平面 \(A_1AMN\perp \text{平面 }EB_1C_1F\)；
\item 设 \(O\) 为 \(\triangle A_1B_1C_1\) 的中心. 若 \(AO//\text{平面 }EB_1C_1F\)，且 \(AO=AB\)，求直线 \(B_1E\) 与平面 \(A_1AMN\) 所成角的正弦值.
\end{enumerate}

\begin{center}
\bitmapinclude[max width=6.0cm,max totalheight=3.4cm]{img/2020/national_paper_2_science/q20_fig1.png}
\end{center}
\end{problem}

\begin{solution}
（1）因为 \(M,N\) 分别为 \(BC,B_1C_1\) 的中点，且底面为正三角形，所以 \(B_1N=BM\)，四边形 \(BB_1NM\) 为矩形，且 \(A_1N\perp B_1C_1\). 于是 \(BB_1//MN\). 又因为 \(AA_1//BB_1\)，所以 \(AA_1//MN\). 又因为 \(MN\perp B_1C_1\)，\(A_1N\perp B_1C_1\)，\(MN\cap A_1N=N\)，所以 \(B_1C_1\perp \text{平面 }A_1AMN\). 而 \(B_1C_1\subset \text{平面 }EB_1C_1F\)，故 \(\text{平面 }A_1AMN\perp \text{平面 }EB_1C_1F\).

（2）因为三棱柱上下底面平行，平面 \(EB_1C_1F\) 与上下底面分别交于 \(B_1C_1,EF\)，所以 \(EF//B_1C_1//BC\). 又因为 \(AO//\text{平面 }EB_1C_1F\)，\(AO\subset \text{平面 }AMNA_1\)，\(\text{平面 }AMNA_1\cap \text{平面 }EB_1C_1F=PN\)，所以 \(AO//PN\). 从而四边形 \(APNO\) 为平行四边形. 因为 \(O\) 是正三角形的中心，且 \(AO=AB\)，所以 \(A_1N=3ON\)，\(AM=3AP\)，\(PN=BC=B_1C_1=3AP=3EF\). 由（1）知直线 \(B_1E\) 在平面 \(A_1AMN\) 内的投影为 \(PN\)，因此直线 \(B_1E\) 与平面 \(A_1AMN\) 所成角即为等腰梯形 \(EFC_1B_1\) 中 \(B_1E\) 与 \(PN\) 所成角. 在等腰梯形 \(EFC_1B_1\) 中，令 \(EF=1\). 过 \(E\) 作 \(EH\perp B_1C_1\) 于 \(H\)，则 \(PN=B_1C_1=EH=3\)，\(B_1H=1\). 所以 \(B_1E=\sqrt{3^2+1^2}=\sqrt{10}\). 于是 \(\sin \angle B_1EH=\frac{B_1H}{B_1E}=\frac{1}{\sqrt{10}}=\frac{\sqrt{10}}{10}\). 故直线 \(B_1E\) 与平面 \(A_1AMN\) 所成角的正弦值为 \(\frac{\sqrt{10}}{10}\).
\end{solution}

\begin{problem}
已知函数 \(f(x)=\sin^2 x\sin 2x\).

\begin{enumerate}
\item 讨论 \(f(x)\) 在区间 \((0,\pi)\) 的单调性；
\item 证明：\(|f(x)|\le \dfrac{3\sqrt3}{8}\)；
\item 设 \(n\in\N^*\)，证明：\(\sin^2 x\sin^2 2x\sin^2 4x\cdots \sin^2 2^n x\le \dfrac{3^n}{4^n}\).
\end{enumerate}

\end{problem}

\begin{solution}
（1）因为
\(f(x)=\sin^2 x\sin 2x=2\sin^3 x\cos x\)，所以
\examdisplayaligned{
f'(x)
&=2\sin^2 x\left( 3\cos^2 x-\sin^2 x \right) \\
&=2\sin^2 x\left( 3-4\sin^2 x \right) \\
&=2\sin^2 x\left[ 3-2\left( 1-\cos 2x \right) \right] \\
&=2\sin^2 x\left( 1+2\cos 2x \right).
}

令 \(f'(x)=0\)，解得 \(x=\frac{\pi}{3}\)或\(x=\frac{2\pi}{3}\). 当 \(x\in\left( 0,\frac{\pi}{3} \right)\) 或 \(x\in\left( \frac{2\pi}{3},\pi \right)\) 时，\(f'(x)>0\)；当 \(x\in\left( \frac{\pi}{3},\frac{2\pi}{3} \right)\) 时，\(f'(x)<0\). 因此 \(f(x)\) 在 \(\left( 0,\frac{\pi}{3} \right)\)，\(\left( \frac{2\pi}{3},\pi \right)\) 上单调递增，在 \(\left( \frac{\pi}{3},\frac{2\pi}{3} \right)\) 上单调递减.

（2）因为 \(f(0)=f(\pi)=0\)，由（1）可知 \(f\left( \frac{\pi}{3} \right)=\frac{3\sqrt3}{8}\)，\(f\left( \frac{2\pi}{3} \right)=-\frac{3\sqrt3}{8}\). 所以 \(f(x)_{\max}=\frac{3\sqrt3}{8}\)，\(f(x)_{\min}=-\frac{3\sqrt3}{8}\). 又因为 \(f(x)\) 为周期函数，所以 \(|f(x)|\le \frac{3\sqrt3}{8}\).

（3）由（2）可知，对任意实数 \(t\)，都有 \(\left| \sin^2 t\sin 2t \right|\le \frac{3\sqrt3}{8}=\left( \frac{3}{4} \right)^{\frac{3}{2}}\). 令 \(a_k=\left| \sin 2^k x \right|\ (k=0,1,\cdots,n)\)，则对 \(k=0,1,\cdots,n-1\)，都有 \(a_k^2a_{k+1}\le \left( \frac{3}{4} \right)^{\frac{3}{2}}\). 将上式两边同时取 \(\dfrac{2}{3}\) 次方，得 \(a_k^{\frac{4}{3}}a_{k+1}^{\frac{2}{3}}\le \frac{3}{4}\). 将上述 \(n\) 个不等式相乘，得 \(a_0^{\frac{4}{3}}a_1^2a_2^2\cdots a_{n-1}^2a_n^{\frac{2}{3}}\le \frac{3^n}{4^n}\). 又因为 \(0\le a_0\le 1\)，\(0\le a_n\le 1\)，所以 \(a_0^2a_1^2a_2^2\cdots a_{n-1}^2a_n^2\le a_0^{\frac{4}{3}}a_1^2a_2^2\cdots a_{n-1}^2a_n^{\frac{2}{3}}\). 因此 \(\sin^2 x\sin^2 2x\sin^2 4x\cdots \sin^2 2^n x=a_0^2a_1^2a_2^2\cdots a_{n-1}^2a_n^2\le \frac{3^n}{4^n}\).
\end{solution}

\begin{problem}
已知曲线 \(C_1,C_2\) 的参数方程分别为 \(C_1: x=4\cos^2\theta\)，\(y=4\sin^2\theta\)（\(\theta\) 为参数），\(C_2: x=t+\dfrac{1}{t}\)，\(y=t-\dfrac{1}{t}\)（\(t\) 为参数）.

\begin{enumerate}
\item 将 \(C_1,C_2\) 的参数方程化为普通方程；
\item 以坐标原点为极点，\(x\) 轴正半轴为极轴建立极坐标系. 设 \(C_1,C_2\) 的交点为 \(P\)，求圆心在极轴上，且经过极点和 \(P\) 的圆的极坐标方程.
\end{enumerate}

\end{problem}

\begin{solution}
（1）对于曲线 \(C_1\)，因为 \(x=4\cos^2\theta\)，\(y=4\sin^2\theta\)，所以 \(x+y=4\)，\(0\le x\le 4\). 故 \(C_1\) 的普通方程为 \(x+y=4\)（且 \(0\le x\le 4\)）.

对于曲线 \(C_2\)，因为 \(x=t+\frac{1}{t}\)，\(y=t-\frac{1}{t}\)，所以 \(x^2-y^2=\left( t+\frac{1}{t} \right)^2-\left( t-\frac{1}{t} \right)^2=4\). 故 \(C_2\) 的普通方程为 \(x^2-y^2=4\).

（2）法一：由 \(x+y=4\) 和 \(\dfrac{x^2}{4}-\dfrac{y^2}{4}=1\) 解得 \(x=\frac{5}{2}\)，\(y=\frac{3}{2}\)，因此 \(P\left( \frac{5}{2},\frac{3}{2} \right)\). 设所求圆的圆心为 \((x_0,0)\)，则由圆经过极点和点 \(P\)，得 \(x_0^2=\left( x_0-\frac{5}{2} \right)^2+\frac{9}{4}\). 解得 \(x_0=\frac{17}{10}\). 故所求圆的极坐标方程为 \(\rho=\frac{17}{5}\cos\theta\).

法二：由 \(x+y=4\) 和 \(x^2-y^2=4\) 得 \(x-y=1\)，所以 \(x=\frac{5}{2}\)，\(y=\frac{3}{2}\)，即 \(P\left( \frac{5}{2},\frac{3}{2} \right)\). 设圆的方程为 \(\left( x-a \right)^2+y^2=r^2\). 由于圆经过点 \(P\) 和原点，所以 \(a^2=r^2\)，且 \(\left( \dfrac{5}{2}-a \right)^2+\left( \dfrac{3}{2} \right)^2=r^2\). 解得 \(a=\frac{17}{10}\)，\(r^2=\frac{289}{100}\). 因此圆的方程为 \(\left( x-\frac{17}{10} \right)^2+y^2=\frac{289}{100}\)，即 \(x^2+y^2-\frac{17}{5}x=0\). 化为极坐标方程为 \(\rho=\frac{17}{5}\cos\theta\).
\end{solution}

\begin{problem}
已知函数 \(f(x)=|x-a^2|+|x-2a+1|\).

\begin{enumerate}
\item 当 \(a=2\) 时，求不等式 \(f(x)\ge 4\) 的解集；
\item 若 \(f(x)\ge 4\)，求 \(a\) 的取值范围.
\end{enumerate}

\end{problem}

\begin{solution}
（1）当 \(a=2\) 时，\(f(x)=|x-4|+|x-3|\)，所以当 \(x\le 3\) 时，\(f(x)=-2x+7\)；当 \(3<x<4\) 时，\(f(x)=1\)；当 \(x\ge 4\) 时，\(f(x)=2x-7\).

当 \(x\le 3\) 时，不等式 \(f(x)\ge 4\) 化为 \(-2x+7\ge 4\). 解得 \(x\le \frac{3}{2}\).

当 \(3<x<4\) 时，不等式 \(f(x)\ge 4\) 化为 \(1\ge 4\)，无解.

当 \(x\ge 4\) 时，不等式 \(f(x)\ge 4\) 化为 \(2x-7\ge 4\). 解得 \(x\ge \frac{11}{2}\).

所以当 \(a=2\) 时，不等式 \(f(x)\ge 4\) 的解集为 \(\left\{ x\mid x\le \frac{3}{2}\ \text{或}\ x\ge \frac{11}{2} \right\}\).

（2）由绝对值不等式的性质，\(f(x)=|x-a^2|+|x-2a+1|\ge \left| x-a^2-\left( x-2a+1 \right) \right|=|(a-1)^2|=(a-1)^2\). 又因为 \(f(x)\ge 4\)，所以 \((a-1)^2\ge 4\). 解得 \(a-1\le -2\ \text{或}\ a-1\ge 2\)，即 \(a\le -1\ \text{或}\ a\ge 3\). 所以 \(a\) 的取值范围为 \((-\infty,-1]\cup[3,+\infty)\).
\end{solution}
